% Focused framework presentation.
\section[Geometry and transport]{From flow geometry to probability transport}

\begin{frame}[t]{The derivative of a flow}
\hypertarget{main:variational}{}
\heading{Perturbing the initial state}
For a $C^1$ state flow and a small displacement $\delta x$,
\begin{equation*}
  \Phi_{s,t}(x+\delta x)=\Phi_{s,t}(x)+M_{s,t}(x)\delta x+o(\|\delta x\|),
  \qquad M_{s,t}(x)=D_x\Phi_{s,t}(x).
\end{equation*}
\heading{Variational equation}
Differentiating $\partial_t\Phi_{s,t}(x)=v_t(\Phi_{s,t}(x))$ in $x$ gives
\begin{equation*}
  \begin{cases}
    \partial_t M_{s,t}(x)=Dv_t(\Phi_{s,t}(x))M_{s,t}(x),\\
    M_{s,s}(x)=I_d.
  \end{cases}
\end{equation*}
The nonlinear trajectory determines the coefficients of a linear matrix ODE.
\heading{Composition of derivatives}
\begin{equation*}
  M_{s,t}(x)=M_{r,t}(\Phi_{s,r}(x))M_{s,r}(x).
\end{equation*}
The intermediate matrix is evaluated at the transported state.
\slidecite{Manuscript \S\,1.2.1, variational equation; Assessment 9.3--4.}
\end{frame}

\begin{frame}[t]{Time ordering in matrix evolution}
\hypertarget{main:timeordering}{}
Fix a trajectory and set $A_t=Dv_t(\Phi_{s,t}(x))$. Then
\begin{equation*}
  \dot M_t=A_tM_t,\qquad M_s=I_d.
\end{equation*}
\heading{Successive substitution, for $t\geq s$}
\begin{equation*}
  M_t=I_d+\int_s^t A_{r_1}\dd r_1
  +\int_s^t\int_s^{r_1}A_{r_1}A_{r_2}\dd r_2\dd r_1+\cdots.
\end{equation*}
Later times appear on the left in each operator product.
\begin{columns}[T,onlytextwidth]
  \begin{column}{.48\textwidth}
    \heading{General coefficients}
    \begin{equation*}
      M_t=\mathcal T\exp\!\left(\int_s^t A_r\dd r\right).
    \end{equation*}
    The order matters when $A_rA_u\ne A_uA_r$.
  \end{column}
  \begin{column}{.47\textwidth}
    \heading{Pairwise commuting coefficients}
    If $A_rA_u=A_uA_r$ for every pair,
    \begin{equation*}
      M_t=\exp\!\left(\int_s^t A_r\dd r\right).
    \end{equation*}
  \end{column}
\end{columns}
\slidecite{Manuscript \S\,1.2.1, Peano--Baker series, and Appendix A.5.}
\end{frame}

\begin{frame}[t]{Jacobian determinant and local volume}
\hypertarget{main:jacobian}{}
\heading{Liouville--Jacobi formula}
For a $C^1$ diffeomorphic flow, define $J_{s,t}(x)=\det M_{s,t}(x)$. Then
\begin{equation*}
  \partial_t J_{s,t}(x)
  =\nabla\cdot v_t(\Phi_{s,t}(x))J_{s,t}(x),\qquad J_{s,s}(x)=1,
\end{equation*}
and therefore
\begin{equation*}
  \boxed{J_{s,t}(x)=\exp\!\left(\int_s^t\nabla\cdot v_r(\Phi_{s,r}(x))\dd r\right)>0.}
\end{equation*}
\heading{Geometric meaning}
\begin{itemize}\setlength{\itemsep}{7pt}
  \item $J_{s,t}$ is the local volume ratio of a transported material element.
  \item Positive divergence gives instantaneous expansion. Negative divergence gives contraction.
  \item A divergence-free field gives $J_{s,t}=1$, allowing shear and deformation.
\end{itemize}
\slidecite{Manuscript \S\,1.2.1; Assessment 9.3--4. \hyperlink{app:jacobian}{Jacobi's formula: Appendix A8.}}
\end{frame}

\begin{frame}[t]{Volume change in the nonlinear example}
\hypertarget{main:nonlinear-jacobian}{}
For $\dot x=0.2x$, $\dot y=0.1y+0.3x^2$, set $\tau=t-s$. The exact flow is
\begin{equation*}
  \Phi_{s,t}(x,y)=
  \begin{pmatrix}e^{0.2\tau}x\\ e^{0.1\tau}y+x^2(e^{0.4\tau}-e^{0.1\tau})\end{pmatrix}.
\end{equation*}
\heading{Differentiate the explicit solution}
\begin{equation*}
  M_{s,t}(x,y)=
  \begin{pmatrix}
    e^{0.2\tau}&0\\
    2x(e^{0.4\tau}-e^{0.1\tau})&e^{0.1\tau}
  \end{pmatrix},\qquad J_{s,t}=e^{0.3\tau}.
\end{equation*}
\heading{Check against divergence}
\begin{equation*}
  \nabla\cdot v=0.2+0.1=0.3,\qquad
  \exp\!\left(\int_s^t0.3\dd r\right)=e^{0.3\tau}.
\end{equation*}
The quadratic term changes shape while the area factor is independent of the initial state.
The explicit solution guarantees completeness even though the field is not globally Lipschitz.
\slidecite{Manuscript \S\,1.2.1, nonlinear volume-deformation example. \hyperlink{main:animation}{Return to the animation.}}
\end{frame}

\begin{frame}[t]{A nonlinear flow in motion}
\hypertarget{main:animation}{}
\begin{equation*}
  R=[-1,1]\times[-0.7,0.7],\qquad 0\leq t\leq1.5.
\end{equation*}
\begin{center}
  \animategraphics[
    width=\linewidth,
    autoplay,
    autopause,
    loop,
    controls,
    poster=first,
    label=nonlinear-flow
  ]{20}{\flowframeprefix}{0}{60}
\end{center}
\smallskip
The same flow moves every grid point. The quadratic coupling bends the rectangle.
\slidecite{Manuscript Example 1.15 and its accompanying 61-frame animation.}
\end{frame}

\begin{frame}[t]{Density transport by a diffeomorphism}
\hypertarget{main:density}{}
Assume $\mu_0(\dd x)=\rho_0(x)\dd x$ and a $C^1$ diffeomorphic flow.
\heading{Change of variables}
\begin{equation*}
  \rho_t(y)=\rho_0(\Phi_{0,t}^{-1}(y))\left|\det D\Phi_{0,t}^{-1}(y)\right|.
\end{equation*}
Along $x_t=\Phi_{0,t}(x_0)$, the formula becomes
\begin{equation*}
  \boxed{\rho_t(x_t)J_{0,t}(x_0)=\rho_0(x_0).}
\end{equation*}
\heading{The divergence controls density along a trajectory}
\begin{equation*}
  \rho_t(x_t)=\rho_0(x_0)
  \exp\!\left(-\int_0^t\nabla\cdot v_r(x_r)\dd r\right).
\end{equation*}
Expansion lowers density within a material element. Its probability mass is conserved:
\begin{equation*}
  \mu_t(\Phi_{0,t}(A))=\mu_0(A).
\end{equation*}
These density identities hold almost everywhere for $\rho_0\in L^1$.
\slidecite{Manuscript \S\,1.2.2; Assessment 9.7--8. \hyperlink{app:density}{Change of variables: Appendix A10.}}
\end{frame}

\begin{frame}[t]{The classical continuity equation}
\hypertarget{main:continuity}{}\hypertarget{main:weak}{}
For the transport representative of the density, differentiation along almost every trajectory gives
\begin{equation*}
  \frac{\dd}{\dd t}\rho_t(x_t)=-\rho_t(x_t)\nabla\cdot v_t(x_t).
\end{equation*}
If $\rho(t,x)$ is also $C^1$, the chain rule yields
\begin{equation*}
  \partial_t\rho_t+v_t\cdot\nabla\rho_t=-\rho_t\nabla\cdot v_t,
  \qquad \boxed{\partial_t\rho_t+\nabla\cdot(\rho_tv_t)=0.}
\end{equation*}
\heading{Generator and adjoint}
For $\mathcal L_t f=v_t\cdot\nabla f$ and compactly supported smooth $f$,
\begin{equation*}
  \int(\mathcal L_t f)\rho\dd x
  =-\int f\nabla\cdot(v_t\rho)\dd x,
  \qquad \mathcal L_t^*\rho=-\nabla\cdot(v_t\rho).
\end{equation*}
Thus $\partial_t\rho_t=\mathcal L_t^*\rho_t$.
\heading{Measure formulation}
Without a density, use $\partial_t\mu_t+\nabla\cdot(v_t\mu_t)=0$ in distributions,
under local flux integrability.
\slidecite{Manuscript \S\,1.2.2. \hyperlink{app:weak}{Weak formulation: Appendix A11.}}
\end{frame}
